% Generated by scripts/export_feature_tables.py -- do not edit by hand.
\begin{description}
  \item[\texttt{user\_skin\_type}] Core demographic attribute driving taste. Known at signup, so it is available for every user including ones with zero training history -- the attribute a brand-new user supplies is exactly what lets the hybrid model place them near users with similar skin.
  \item[\texttt{user\_skin\_tone}] Secondary demographic attribute (shade matching for base/coverage products). Same availability argument as skin type.
  \item[\texttt{user\_identity}] Per-user ID feature, LightFM-style: lets the model retain personalised signal for warm users without it being the *only* signal. The vocabulary is restricted to users seen in training, so a cold user genuinely has no ID token available -- this is the feature that operationalizes "cold start" rather than merely asserting it.
  \item[\texttt{item\_brand}] Brand loyalty is a strong prior independent of any single product's review history, and is known the moment a product is catalogued.
  \item[\texttt{item\_category}] Coarse product type (serum vs. cleanser vs. sunscreen). Known at catalogue time, safe for cold items by construction.
  \item[\texttt{item\_price\_bucket}] Discretized price tier. Budget shoppers and prestige shoppers cluster on this signal regardless of brand or category.
  \item[\texttt{item\_ingredient\_tags}] The content signal that actually explains taste: a Dry-skin user's affinity is for Hyaluronic-Acid / Ceramide products specifically, not for "Moisturizers" as a category. This is what lets the model recommend a brand-new product that happens to share tags with liked products.
  \item[\texttt{item\_identity}] Mirrors user\_identity: vocabulary restricted to items seen in training, so a cold item has no ID token and must be scored through its brand/category/ingredient tokens alone.
  \item[\texttt{user\_review\_count\_bucket}] Confidence weighting: a user with 40 training reviews has a much more reliable taste signal than one with 3. Computed with a strict `<` cutoff against the user's own interaction\_order so it can never see the very review it is being used to help predict.
  \item[\texttt{item\_popularity\_bucket}] Popularity prior used by the most-popular baseline and as part of the cold-start fallback ordering. Fitted on the training partition only, so held-out interactions never inflate a product's apparent popularity.
  \item[\texttt{item\_catalogue\_context}] Distinguishes an "established" item (enough training reviews to trust its popularity/ingredient signal) from a "new\_arrival" (sparse or zero training reviews). Cold items always fall back to this value, documenting that they are indistinguishable from genuinely new catalog entries from the model's point of view.
\end{description}
